\section{Three statements}
\label{sec:ch04-three-statements}
\index{DataLoader|(}

The pattern is older than GraphQL and it is named after the library that
popularised it. DataLoader's README credits Nick Schrock with the original
Loader API, written at Facebook in 2010 to coalesce a collection of key-value
back ends, and describes what it does in one sentence: batching and caching over
remote data sources~\autocite{schrock2026dataloader}. Both halves matter here,
and one of them is doing more work than people expect.

\index{DataLoader!source generator}%
In HotChocolate 16 you do not write the class. You write a static method that
takes a list of keys and returns a keyed result, mark it, and the source
generator writes the DataLoader and its interface at build
time~\autocite{chillicream2026dataloaderdocs}:

\begin{minted}{csharp}
public static class AccountDataLoaders
{
    [DataLoader]
    public static Task<IReadOnlyDictionary<Guid, Customer>> GetCustomerByIdAsync(
        IReadOnlyList<Guid> ids,
        AccountsService accounts,
        CancellationToken cancellationToken)
        => accounts.GetCustomersByIdsAsync(ids, cancellationToken);
}
\end{minted}

That produces \code{CustomerByIdDataLoader} and \code{ICustomerByIdDataLoader}.
The name is derived by stripping a leading \code{Get}, then a trailing
\code{Async}, then a trailing \code{DataLoader}, and appending \code{DataLoader}
to what is left. It is worth knowing the rule rather than guessing at the result,
because the interface name is what you type in the resolver.

\index{DataLoader!partial classes}%
Two smaller things about that snippet. The class is not \code{partial}, and does
not need to be, which sets it apart from every other generated thing in Mosaic:
\code{[QueryType]} and \code{[ObjectType<T>]} both extend the class they are on,
and \code{[DataLoader]} generates a separate class beside it. And there is no
registration anywhere. The documentation shows an assembly attribute,
\code{DataLoaderModule}, for that~\autocite{chillicream2026dataloaderdocs}. But
Mosaic already declares the type module attribute
chapter~\ref{ch:hotchocolate-quickly} introduced:

\begin{minted}{csharp}
[assembly: Module("Mosaic")]
\end{minted}

The type module generator emits an \code{AddDataLoader} call for every generated
loader into the same \code{AddMosaic()} it already emits the types into. One
attribute, declared for another reason two chapters ago, covers both.

\subsection*{Three shapes, and one of them is a trap}

\index{DataLoader!kinds}%
The generator picks what kind of DataLoader to build from the return type, in
this order: a dictionary keyed by the key type is a \emph{batch} loader, an
\code{ILookup} is a \emph{group} loader, and anything else is treated as a
\emph{cache} loader over a single key.

The distinction that matters in practice is not the mechanism. All three batch;
all three cache. What differs is what happens to a key that has no answer.

\begin{minted}{csharp}
[DataLoader]
public static Task<ILookup<Guid, Review>> GetReviewsByProductIdAsync(
    IReadOnlyList<Guid> productIds,
    ReviewsService reviews,
    CancellationToken cancellationToken)
    => reviews.GetReviewsByProductIdsAsync(productIds, cancellationToken);
\end{minted}

\code{ILookup} answers an unknown key with an empty sequence. Three of Mosaic's
twenty-five products have never been reviewed, and the \code{reviews} field is a
non-nullable list, so this return type is the difference between a correct
response and three null violations. Had I written
\code{Dictionary<Guid, Review[]>} instead, the code would have compiled, the
batching would have worked identically, and the three unreviewed products would
have come back as errors. The documentation is explicit that a dictionary whose
value happens to be an array is still a batch loader, and that a missing key in
a batch loader is null~\autocite{chillicream2026dataloaderdocs}. It is the kind
of thing you read past until it bites.

\subsection*{The resolvers, before and after}

\index{DataLoader!in resolvers}%
Here is \code{Review.author} as chapter~\ref{ch:hotchocolate-quickly} left it,
and as this chapter leaves it. The declaration is the only thing that changed:

\begin{minted}[fontsize=\footnotesize]{csharp}
public static async Task<Customer> GetAuthorAsync(
    [Parent] Review review,
    AccountsService accounts,
    CancellationToken cancellationToken)
{
    var author = await accounts.GetCustomerByIdAsync(review.CustomerId, cancellationToken);
\end{minted}

\begin{minted}[fontsize=\footnotesize]{csharp}
public static async Task<Customer> GetAuthorAsync(
    [Parent] Review review,
    ICustomerByIdDataLoader customerById,
    CancellationToken cancellationToken)
{
    var author = await customerById.LoadAsync(review.CustomerId, cancellationToken);
\end{minted}

One parameter type, one method name. The field keeps its name, its type and its
nullability, so the schema does not move. That property is why this fix can be
applied to a live service one resolver at a time, and why nobody has to be told
it is happening.

\index{DataLoader!request scope}%
The same interface appears in \code{Order.customer}, and that is not tidiness. A
DataLoader instance and its cache live for exactly one
request~\autocite{chillicream2026dataloaderdocs}, so a query that reaches a
customer through their orders and again through their reviews fetches them once
and hands both fields the same object. Consistency inside a response is a
property you get for free here, and it is worth naming because it is the sort of
thing that is very expensive to add later.

\subsection*{The number}

\index{N+1 problem!fixed}%
Same query, same seed data, same machine:

\begin{minted}[fontsize=\footnotesize]{text}
parse - validate - compile - coerce - execute 4.365ms total 4.414ms
    (document cache hit, operation cache hit, 146 resolvers, 3 SQL)
\end{minted}

A hundred and forty-six resolvers, and three statements. Nine of ten warm runs
came in between 3.9 and 5.8~ms, against 9.2 to 14.2 for the same query at the
previous tag, and I measured the two back to back in one sitting because
otherwise the comparison is between two different machine loads rather than
between two versions of Mosaic. The caveat from
section~\ref{sec:ch04-what-a-lookup-costs} still applies: the database is on the
same machine, so the ratio is what travels rather than the milliseconds.

The statements themselves are worth reading, because they are ordinary:

\begin{minted}[fontsize=\footnotesize]{sql}
SELECT p.id, p.category, p.description, p.sku, p.title
FROM products AS p
ORDER BY p.id

SELECT r.id, r.body, r.created_at, r.customer_id, r.product_id, r.rating
FROM reviews AS r
WHERE r.product_id = ANY (@productIds)
ORDER BY r.created_at, r.id

SELECT c.id, c.display_name, c.email
FROM customers AS c
WHERE c.id = ANY (@ids)
\end{minted}

\index{PostgreSQL!ANY versus IN}%
\code{Contains} over a list becomes \code{= ANY (@keys)} on Npgsql rather than
the \code{IN (@p0, @p1, ...)} that other providers emit. That is one parameter
instead of one per key, which means one query plan whatever the page size. A
provider that expands the list produces a differently shaped statement for every
distinct key count, and a plan cache full of near-identical entries. If you take
one thing from this listing that is not about GraphQL, take that.

Figure~\ref{fig:ch04-batching} draws the same request twice, once for each tag.

\begin{figure}[htbp]
  \centering
  % The same 146 resolvers, with and without DataLoaders. Referenced from
% chapter 04. Bare tikzpicture: the figure environment, caption and label live
% at the call site. Uses only arrows.meta and positioning.
\begin{tikzpicture}[
    font=\small,
    lvl/.style={draw, rounded corners=2pt, minimum width=30mm,
                minimum height=7mm, align=center, font=\scriptsize},
    db/.style={draw, thick, rounded corners=2pt, minimum width=30mm,
               minimum height=8mm, align=center, font=\scriptsize},
    batch/.style={draw, fill=black!8, rounded corners=2pt, minimum width=30mm,
                  minimum height=7mm, align=center, font=\scriptsize},
    thin arrow/.style={-{Stealth[length=1.6mm]}, gray},
    fat/.style={-{Stealth[length=2mm]}, thick},
    note/.style={font=\scriptsize\itshape, align=center},
    head/.style={font=\small\bfseries}
  ]

  % ---- left: one resolver, one statement -------------------------------
  \node[head] at (-3.4,1.1) {tag \texttt{ch04-ef}};

  \node[lvl] (lp) at (-3.4,0.2)  {\texttt{products}\\1 resolver};
  \node[lvl] (lr) at (-3.4,-1.6) {\texttt{reviews}\\25 resolvers};
  \node[lvl] (la) at (-3.4,-3.4) {\texttt{author}\\120 resolvers};
  \node[db]  (ldb) at (-3.4,-5.3) {PostgreSQL};

  \draw[fat] (lp) -- (lr);
  \draw[fat] (lr) -- (la);

  % Many thin arrows: every resolver goes on its own.
  \foreach \x in {-4.4,-4.0,...,-2.4}{%
    \draw[thin arrow] (\x,-3.75) -- (\x,-4.9);
  }
  \foreach \x in {-4.2,-3.8,...,-2.6}{%
    \draw[thin arrow] (\x,-1.95) .. controls (\x-0.3,-3.4) .. (\x-0.15,-4.9);
  }
  \draw[thin arrow] (-3.4,-0.15) .. controls (-5.4,-2.5) .. (-4.4,-4.9);

  \node[note] at (-3.4,-6.1) {146 statements\\one per question};

  % ---- right: keys first, then three statements ------------------------
  \node[head] at (3.4,1.1) {tag \texttt{ch04}};

  \node[lvl] (rp) at (3.4,0.2)  {\texttt{products}\\1 resolver};
  \node[lvl] (rr) at (3.4,-1.6) {\texttt{reviews}\\25 resolvers};
  \node[lvl] (ra) at (3.4,-3.4) {\texttt{author}\\120 resolvers};

  \node[batch] (b1) at (3.4,-4.5) {batch: 12 keys};
  \node[db]    (rdb) at (3.4,-5.9) {PostgreSQL};

  \draw[fat] (rp) -- (rr);
  \draw[fat] (rr) -- (ra);

  % The 120 resolvers hand keys to one batch.
  \foreach \x in {2.4,2.8,...,4.4}{%
    \draw[thin arrow] (\x,-3.75) -- (\x,-4.2);
  }
  \draw[fat] (b1) -- (rdb);

  \node[note] at (3.4,-6.7) {3 statements\\one per batch};

  \node[note, anchor=west, text width=32mm] at (5.1,-4.5)
    {120 calls, 12 keys:\\the cache answers a\\repeated key before\\a batch sees it};

\end{tikzpicture}

  \caption{The same hundred and forty-six resolvers, twice. On the left every
  resolver reaches the database on its own account. On the right the keys are
  gathered before anything is fetched, and the hundred and twenty author
  resolvers collapse to twelve keys rather than a hundred and twenty, because
  deduplication happens in the cache before a key ever reaches a batch.}
  \label{fig:ch04-batching}
\end{figure}

\subsection*{The half that is not batching}

\index{DataLoader!deduplication}%
The third statement carries twelve keys. The query asked for the authors of a
hundred and twenty reviews, and Mosaic has twelve customers.

This is the part of DataLoader that gets less attention than the batching and
that does more work in a real graph. \code{LoadAsync} consults a promise cache
\emph{before} it touches a batch: the key is looked up, and a repeated key gets
the promise the first caller made rather than a place in the queue. So the
hundred and twenty calls produce twelve entries, and the fetch method is handed
twelve keys. Batching turned a hundred and twenty statements into one;
deduplication decided what that one statement had to ask for.

The two work together and they fail together. A cache key that varies when it
should not, an identifier that arrives as a string in one place and a
\code{Guid} in another, and the batch is suddenly a hundred and twenty keys
long. It still works, it is still one statement, and it is doing ten times the
work with nothing in the response to say so. Which is a good argument for
keeping a count of statements per request somewhere you will see it.

Which leaves the question the numbers so far have not answered: what decides
that the batch is finished?
\index{DataLoader|)}
